\section{The Proportionality Corridor}
\label{sec:corridor}

\subsection{Definition}

We define the \emph{proportionality corridor} as the set of seat outcomes
within one seat of strict proportional representation.  For a state with
$k$ seats and a statewide Democratic vote share $v_D$, the proportional
seat count is $\lfloor k \cdot v_D \rceil$ (rounded to nearest integer).
The corridor is:

\begin{equation}
  C_{\text{prop}} = \{S_D : |S_D - \lfloor k \cdot v_D \rceil| \leq 1\}.
  \label{eq:corridor}
\end{equation}

The \emph{corridor fraction} is the fraction of ensemble plans whose partisan
outcome falls in $C_{\text{prop}}$:

\begin{equation}
  \phi = \frac{1}{N} \#\{P \in \mathcal{E} : S_D(P) \in C_{\text{prop}}\}.
  \label{eq:corridor-fraction}
\end{equation}

\subsection{Corridor Fractions Across States}

For competitive states, $\phi$ is typically 20--40\%: most valid plans fall
within one seat of proportional, but many do not.  For heavily sorted states,
$\phi$ can be much smaller.

\begin{center}
\begin{tabular}{lcccc}
\toprule
State & $v_D$ (2020) & Proportional seats & Corridor & $\phi$ \\
\midrule
NC & 0.491 & 7     & $\{6, 7, 8\}$ & 42\% \\
WI & 0.497 & 4     & $\{3, 4, 5\}$ & 61\% \\
GA & 0.464 & 6.5   & $\{6, 7\}$   & 28\% \\
PA & 0.502 & 8.5   & $\{8, 9\}$   & 55\% \\
MN & 0.531 & 4.2   & $\{4, 5\}$   & 59\% \\
\bottomrule
\end{tabular}
\end{center}

\subsection{AR's Position in the Corridor}

The \ar{} plan falls inside the proportionality corridor for NC (7D = proportional),
WI (4D = proportional), and PA (8D, within one of proportional).  It falls
outside the corridor for Georgia (5D vs.\ proportional of 6--7D).

\begin{keybox}
\textbf{Key finding:} \ar{} falls inside the proportionality corridor for
5 of 6 studied states (including all competitive states).  The single
exception, Georgia, is explained by geographic sorting under the corridor's
definition, not by algorithmic bias.
\end{keybox}

\subsection{Minnesota: An Additional Competitive State}

We include Minnesota ($k = 8$, $v_D = 0.531$) as an additional competitive
state not studied in G.1.  The proportional seat count is approximately 4.2,
so the corridor is $\{4, 5\}$.  The \ar{} plan for Minnesota produces 4D/4R,
inside the corridor at the 44th percentile of Democratic seats.  This
confirms the robustness of the finding across competitive states.

\subsection{Interpretation}

Falling inside the proportionality corridor does not mean the plan is
proportional by design.  It means the plan's partisan outcome is within
the range that proportional representation would imply.  This is a weaker
claim---and a more honest one.  The \ar{} algorithm makes no proportionality
guarantee; it guarantees minimum edge cut.  The corridor result shows that,
empirically, minimum edge cut and geographic proportionality are compatible
for competitive states.

For sorted states like Georgia, they are not fully compatible.  This is an
important limitation to acknowledge: the \ar{} algorithm cannot overcome
the geographic sorting of voters.
